% TeX root = ../Main.tex

% First argument to \section is the title that will go in the table of contents. Second argument is the title that will be printed on the page.
\section[Question -- {\it Estimating optical flow}]{}

\subsection{Estimating optical flow}
To estimate optical flow, the Lk algorithm is employed. It is based on three main absumption:
\begin{itemize}
    \item Brightness constancy: $I(x, y, t) = I(x+u \delta t, y+v \delta t, t+\delta t)$
    \item Small motion
    \item Spatial coherence
\end{itemize}

Starting from the brightness constancy, by applying Taylor approximation, we can derive the optical flow constraint equation
$I_x u + I_y v+ I_t = 0$. 
Since this single-pixel equation is underconstrained (1 equation for 2 unknowns), Lucas-Kanade introduces the spatial coherence assumption, yielding $N^2$ equations for an $N \times N$ patch.
$Ax = b$ is so formed as an overdetermined system solvable by Least Squares. The Solution is the eigenvector corresponding to 
the smallest eigenvalue.
This is possible only if $A^T A$ is invertible and well conditioned. Indeed, by the $A^T A$ eigenvalues we can understand
the optical flow, as for the harris corner detector:
\begin{itemize}
    \item $\lambda_1 \approx \lambda_2 \approx 0$ (Flat region): No reliable optical flow can be estimated.
    \item $\lambda_1 \ \lambda_2 \approx 0$ (Edge): Aperture problem; only the normal flow (perpendicular to the edge) can be estimated.
    \item $\lambda_1 \approx \lambda_2 \gg 0$ (Corner / Texture): Optical flow can be estimated reliably in all directions.
\end{itemize}

The **aperture problem** occurs in the edge case, stating that observing a moving edge through a small local aperture 
only allows estimating the motion component perpendicular to the edge, while the parallel (tangent) component 
remains completely ambiguous. Lucas-Kanade mitigates this by pooling information from a spatial neighborhood (patches) 
rather than relying on isolated pixels.